\section{Receiving the word and learning to live by it}

Jesus compares the kingdom of heaven to a very small seed and to leaven hidden in flour. Paul describes people who have received the Gospel amid considerable suffering and have become an example to others. In each case an apparently unpromising beginning produces something that reaches beyond itself. Yet the prayers of this Mass ask for more than growth. They ask that thought become acceptable speech and action, that the offering cleanse and renew those who bring it, and that people already fed with heavenly gifts continue to desire them.

These requests raise two connected questions. How can people exposed to opposition communicate the faith they have received? And what must change within them if that faith is to direct their lives? The first interpretation follows Chrysostom's explanation of the seed and leaven alongside Aquinas's account of the Thessalonians' example. The second follows Jerome's explanation of the Gospel's effect on the whole person alongside Augustine's exposition of prayer from the depths. Their answers meet in ordinary Christian conduct, but they approach its causes differently. One explains how a community can become fruitful under trial; the other explains how understanding and desire can be converted.

The Resumed Sixth Sunday after the Epiphany is celebrated on 15 November 2026 under the universal calendar of the 1962 Missal. It retains the Sixth Sunday's prayers, Epistle and Gospel, with the chants appointed for resumed use after Pentecost. The Preface of the Most Holy Trinity is expressly appointed.\footnote{\textit{Missale Romanum} (1962), pp.~421--422, nos.~1759--1768; Preface p.~293, no.~1082; \textit{Rubricae generales} 18(b); \textit{Rubricae generales Missalis romani} 298.} The following table follows their order in the Mass. Each text is given in the next section, before the detailed commentary and the two interpretations.

\begin{maptable}
Introit & Jeremiah 29:11,12,14; Psalm 84(85):2 & God promises peace and gathering; the psalm verse recalls his favor toward his people.\\
Collect & \textit{Praesta, quaesumus} & Reflection is to issue in words and deeds pleasing to God.\\
Epistle & 1 Thessalonians 1:2--10 & Faith received amid affliction becomes an example, while believers await the risen Son.\\
Gradual & Psalm 43(44):8--9 & Deliverance is attributed to God, whose name will be praised for ever.\\
Alleluia & Psalm 129(130):1--2 & A voice calls to God from the depths.\\
Gospel & Matthew 13:31--35 & Jesus teaches through mustard seed and leaven, and Matthew cites the psalmist's parables.\\
Offertory & Psalm 129(130):1--2 & The petition for a hearing returns, with the opening cry repeated.\\
Secret & \textit{Haec nos oblatio} & The offering is asked to cleanse, renew, govern and protect the people.\\
Preface & Most Holy Trinity & Thanksgiving confesses one God, the distinction of persons, unity of essence and equality of majesty.\\
Communion & Mark 11:24 & Christ assures those who pray of God's hearing.\\
Postcommunion & \textit{Caelestibus, Domine} & Those already nourished ask for continuing desire for the gifts by which they live.\\
\end{maptable}

The four senses developed after each interpretation distinguish what the texts say in their immediate contexts, their fulfillment in Christ and his body, their bearing on conduct, and their final fulfillment. That final perspective is called the anagogical sense. Historical settings and the sourced alternatives for biblical dates are gathered in the Scriptural Date and Location appendix (\ref{sec:date-location}).
